\section{Manifesto Results: Main Application}\label{sec:min-manifesto-results}

The Manifesto application asks whether a learned compression tree can preserve
a social-science measurement target while leaving behind readable intermediate
states. The source documents are party election programs from the Manifesto
Project. The project collects party platforms, releases document text through
the Manifesto Corpus, and releases hand-coded policy-emphasis data through the
Manifesto Data Collection \citep{ManifestoProject2025a,ManifestoCorpus2025}.
Its coding tradition goes back to the Manifesto Research Group and
Comparative Manifestos Project books on spatial policy measurement
\citep{BudgeRobertsonHearl1987,BudgeEtAl2001,KlingemannEtAl2006,VolkensEtAl2013};
\citet{MerzEtAl2016} describe the modern corpus as a multilingual text
resource. The unit of human coding in the Manifesto Project is the
quasi-sentence: a text unit expressing one political statement, assigned to a
policy category. The Benoit benchmark below uses expert means as targets;
quasi-sentence labels are the natural next source of human node labels for C1
and C3 audits.

We use the benchmark organized by \citet{BenoitEtAl2025}. It contains $235$
party platforms matched to expert-survey means on six seven-point policy
dimensions: economic policy, social policy, immigration, European integration,
environmental policy, and decentralization. For each manifesto--dimension
pair, the target is the mean expert placement of the party on that issue
scale. This target differs from both the Manifesto Project quasi-sentence
categories and any individual coder label. A method is evaluated by Pearson
correlation between its root score and the expert mean within each dimension;
the macro score is the unweighted mean of the six dimension correlations.
Pearson is unchanged by the affine rescaling from Benoit's released non-EU
expert means to the $1$--$7$ score scale; corrected MAE and calibration
diagnostics use the normalized $1$--$7$ scale.

Benoit et al. also give the LLM replication target. Their pipeline first
summarizes each manifesto--dimension pair into a short English summary, then
scores that summary on the seven-point expert scale. Their headline ensemble
uses $3$ LLMs, $2$ prompting regimes, and $3$ generated summaries, giving $18$
scores per manifesto--dimension. The three reference values we use are the
split-expert agreement row, macro $r=0.873$; the proprietary LLM ensemble,
macro $r=0.817$; and the matched Gemma-3 open-weight replication, macro
$r=0.793$. The split-expert row is a measurement agreement reference:
experts are split into two groups and the group means are correlated. It
serves as measurement context; no ladder is trained or selected on it.

These results are a root-observed corpus evaluation. Each held-out
manifesto--dimension pair has a full-document expert target, and we evaluate
the tree's root score directly against that target. A node-level local-law
audit is stronger and different: it would sample leaves and merge spans with
logged propensities and compare their local evidence to an application-aligned
oracle. The current section reports the root-observed LLM application and the
available teacher-trace diagnostics, not a completed human node-level audit.

\subsection{Tree Setup and Splits}\label{sec:min-manifesto-tree-setup}

The C-TreePO version keeps Benoit's document-level target but changes the
representation being scored. The query $f$ is the policy score and the
summarizer/merger $g$ is the learned state map. A flat baseline asks $f$ to
score the manifesto text directly. A tree run chunks the manifesto into
contiguous leaves, maps each leaf into a readable policy-evidence state,
recursively merges adjacent states, and applies $f$ at the root. The root
comparison instantiates $f(g(x))\approx f(x)$; each internal merge
instantiates the C3 route
$f(g(g(u)\concat g(v)))\approx f(u\concat v)$.

The current f/g ladder artifacts are a subset of the reference benchmark with
complete text, expert targets, and teacher-trace preprocessing. They contain
$218$ manifestos split into $140$ training documents, $30$ reserved documents,
and $48$ held-out test documents. Each manifesto is materialized as one tree
for each leaf size. The phrase ``$48$ test trees'' therefore means $48$
held-out manifestos, each with one root prediction for a given method and leaf
size. Reported ladder metrics are held-out test metrics. The $30$ reserved
documents are materialized for later checks; the DSPy prompt search described
here uses node examples from the $140$ training documents in its compile loop.
The node count changes with leaf size. In the joint all-six run, $256$-token
leaves create $3{,}249$ total teacher nodes across the $218$ documents, while
$8{,}096$-token leaves create $218$ total teacher nodes. In the
decentralization single-target run, $256$-token leaves create $1{,}841$ total
teacher nodes and $1{,}024$-token leaves create $467$.

The f/g ladder alternates the two operators. First, teacher traces create
target summaries and target scores for every realized node. With the query
fixed, we optimize $g$ so the current query recovers the node target from the
candidate summary. With $g$ fixed, we optimize $f$ on the summaries produced
by the current tree. The table label $f^a g^b$ records the number of
query-side and summarizer-side optimization rounds. The important point is
co-adaptation: the state map is trained against the query, and the query
is trained on the state distribution produced by the map.\footnote{Some
legacy JSON artifacts label the initial query artifact as
\texttt{identity}. Operationally, that row loads the pre-tuned dimension
query, while the initial summarizer is teacher-summary passthrough.}

The result artifacts are generated tables and figures under
\texttt{assets/benoit/}. The headline table is
\texttt{benoit\_comparison\_pearson.tex}; the f/g ladder figures use the
same $140/30/48$ split and report held-out Pearson against expert means.
Character-leaf sweeps and token-leaf prompt ladders are separate axes and are
reported separately below.

\subsection{Six-Dimension Results}\label{sec:min-manifesto-six-dim}

\begin{table}[t]
  \centering
  \resizebox{\linewidth}{!}{% Generated table; do not edit by hand.
% Metric: Pearson r (higher better)
\begin{tabular}{lrrrrrrrr}
\toprule
Method & Econ. & Social & Immig. & EU & Env. & Decent. & Macro & \#/6 \\
\midrule
\multicolumn{9}{l}{\textit{Benoit et al. (2025) reference}} \\
Proprietary ensemble, 18 scores (Fig 1) & $+0.870$ & $+0.920$ & $+0.890$ & $+0.910$ & $+0.820$ & $+0.490$ & $+0.817$ & 6/6 \\
Split-expert reference (Table 3) & $+0.880$ & $+0.910$ & $+0.880$ & $+0.950$ & $+0.840$ & $+0.780$ & $+0.873$ & 6/6 \\
LLaMA-3.3-70B (Table 6) & $+0.840$ & $+0.870$ & $+0.860$ & $+0.860$ & $+0.680$ & $+0.400$ & $+0.752$ & 6/6 \\
DeepSeek-V3 (Table 6) & $+0.840$ & $+0.870$ & $+0.890$ & $+0.860$ & $+0.790$ & $+0.450$ & $+0.783$ & 6/6 \\
Gemma-3-27B-IT (Table 6) & $+0.860$ & $+0.860$ & $+0.890$ & $+0.840$ & $+0.860$ & $+0.450$ & $+0.793$ & 6/6 \\
\midrule
\multicolumn{9}{l}{\textit{Ours: per-dim pipeline $\times$ leaf size (Gemma-4-31B-NVFP4, 1 summarizer + 1 scorer per dim)}} \\
leaf = 64 K chars ($\approx$16K tokens) & $+0.916$ & $+0.817$ & $+0.888$ & $+0.924$ & $+0.854$ & $+0.489$ & $+0.815$ & 6/6 \\
leaf = 32 K chars ($\approx$8K tokens) & $+0.929$ & $+0.857$ & $+0.886$ & $+0.916$ & $+0.877$ & $+0.480$ & $+0.824$ & 6/6 \\
leaf = 24 K chars ($\approx$6K tokens) --- full test n$\approx$215 & $+0.892$ & $+0.840$ & $+0.867$ & $+0.896$ & $+0.814$ & $+0.464$ & $+0.796$ & 6/6 \\
leaf = 16 K chars ($\approx$4K tokens) & $+0.925$ & $+0.847$ & $+0.898$ & $+0.909$ & $+0.831$ & $+0.537$ & $+0.824$ & 6/6 \\
leaf =  8 K chars ($\approx$2K tokens) & $+0.939$ & $+0.869$ & $+0.884$ & $+0.901$ & $+0.839$ & $+0.543$ & $+0.829$ & 6/6 \\
\midrule
\multicolumn{9}{l}{\textit{Ours: combined pipeline $\times$ leaf size (one shared summarizer with union rubric, 6 scores)}} \\
leaf = 64 K chars & $+0.910$ & $+0.867$ & $+0.897$ & $+0.910$ & $+0.795$ & $+0.458$ & $+0.806$ & 6/6 \\
leaf = 32 K chars & $+0.917$ & $+0.849$ & $+0.908$ & $+0.889$ & $+0.749$ & $+0.438$ & $+0.792$ & 6/6 \\
leaf = 24 K chars (full test n=229) & $+0.868$ & $+0.850$ & $+0.880$ & $+0.902$ & $+0.809$ & $+0.396$ & $+0.784$ & 6/6 \\
leaf = 16 K chars & $+0.919$ & $+0.851$ & $+0.916$ & $+0.889$ & $+0.808$ & $+0.441$ & $+0.804$ & 6/6 \\
leaf =  8 K chars & $+0.927$ & $+0.874$ & $+0.911$ & $+0.917$ & $+0.801$ & $+0.413$ & $+0.807$ & 6/6 \\
\midrule
\multicolumn{9}{l}{\textit{Ours: tiny-leaf extensions of the chunk sweep (Gemma-4)}} \\
tree, leaf = 4 K chars & $+0.932$ & $+0.886$ & $+0.885$ & $+0.931$ & $+0.838$ & $+0.580$ & $+0.842$ & 6/6 \\
tree, leaf = 2 K chars & $+0.924$ & $+0.867$ & $+0.887$ & $+0.906$ & $+0.793$ & $+0.584$ & $+0.827$ & 6/6 \\
tree, leaf = 1 K chars (stress test, depth $\geq$6) & $+0.933$ & $+0.859$ & $+0.914$ & $+0.925$ & $+0.838$ & $+0.493$ & $+0.827$ & 6/6 \\
\midrule
\multicolumn{9}{l}{\textit{Ours: concat-no-merge (chunks summarized independently, joined, scored --- tests whether the merge step carries signal)}} \\
concat, leaf = 32 K chars & $+0.929$ & $+0.856$ & $+0.897$ & $+0.920$ & $+0.848$ & $+0.571$ & $+0.837$ & 6/6 \\
concat, leaf = 16 K chars (default) & $+0.931$ & $+0.823$ & $+0.909$ & $+0.912$ & $+0.866$ & $+0.523$ & $+0.827$ & 6/6 \\
concat, leaf =  8 K chars & $+0.932$ & $+0.849$ & $+0.882$ & $+0.915$ & $+0.850$ & $+0.594$ & $+0.837$ & 6/6 \\
\midrule
\multicolumn{9}{l}{\textit{Ours: flat baseline (no chunk, no summary; truncate text and score) --- tests whether tree is needed at all}} \\
flat, truncation = 48 K chars & $+0.936$ & $+0.821$ & $+0.926$ & $+0.889$ & $+0.806$ & $+0.587$ & $+0.827$ & 6/6 \\
flat, truncation = 24 K chars (default) & $+0.929$ & $+0.847$ & $+0.889$ & $+0.882$ & $+0.829$ & $+0.542$ & $+0.820$ & 6/6 \\
flat, truncation = 12 K chars & $+0.929$ & $+0.844$ & $+0.938$ & $+0.880$ & $+0.779$ & $+0.564$ & $+0.822$ & 6/6 \\
flat, truncation =  6 K chars & $+0.926$ & $+0.811$ & $+0.967$ & $+0.881$ & $+0.886$ & $+0.243$ & $+0.786$ & 6/6 \\
\midrule
\multicolumn{9}{l}{\textit{Ours: scorer ablations (Gemma-4, Benoit's GPT-4o summaries held fixed)}} \\
1. Zero-shot scorer & $+0.832$ & $+0.886$ & $+0.858$ & $+0.905$ & $+0.668$ & $+0.311$ & $+0.743$ & 6/6 \\
2. Few-shot optimized scorer & $+0.822$ & $+0.892$ & $+0.853$ & $+0.912$ & $+0.627$ & $+0.297$ & $+0.734$ & 6/6 \\
3. Joint scorer baseline (shared across 6 dims) & $+0.837$ & $+0.881$ & $+0.852$ & $+0.904$ & $+0.658$ & $+0.314$ & $+0.741$ & 6/6 \\
4. Joint scorer, few-shot optimized & $+0.817$ & $+0.891$ & $+0.848$ & $+0.908$ & $+0.700$ & $+0.313$ & $+0.746$ & 6/6 \\
5. Joint scorer, reflective prompt-optimized & $+0.832$ & $+0.882$ & $+0.848$ & $+0.879$ & $+0.688$ & $+0.344$ & $+0.746$ & 6/6 \\
\midrule
\multicolumn{9}{l}{\textit{Ours: exact-model replication (Benoit's Gemma-3-27B-IT BF16)}} \\
Gemma-3-27B scorer-only, OUR scoring rubric & $+0.826$ & $+0.864$ & $+0.846$ & $+0.902$ & $+0.573$ & $+0.329$ & $+0.723$ & 6/6 \\
Gemma-3-27B scorer-only, exact Benoit rubric & $+0.814$ & $+0.895$ & $+0.852$ & $+0.852$ & $+0.616$ & $+0.359$ & $+0.731$ & 6/6 \\
Gemma-3-27B, raw-prompt Benoit format (bare-integer response) & $+0.782$ & $+0.901$ & $+0.854$ & $+0.874$ & $+0.590$ & $+0.368$ & $+0.728$ & 6/6 \\
\bottomrule
\end{tabular}
}
  \caption{\textbf{Pearson $r$ against Benoit expert-survey means.}
  The first block reproduces published Benoit reference rows; the remaining
  blocks report C-TreePO and controls using Gemma-4-31B-NVFP4. The
  per-dimension tree at $8\mathrm{K}$-character leaves reaches macro
  $r=0.829$ while reporting every dimension separately.}
  \label{tab:min-benoit-headline}
\end{table}

Table \ref{tab:min-benoit-headline} gives the full six-dimension benchmark.
The per-dimension tree at $8\mathrm{K}$ characters reaches macro $r=0.829$,
with dimension scores $0.939$ economic, $0.869$ social, $0.884$
immigration, $0.901$ European integration, $0.839$ environment, and $0.543$
decentralization. The corresponding combined-summary tree at the same leaf
size reaches macro $r=0.807$, with $0.927$, $0.874$, $0.911$, $0.917$,
$0.801$, and $0.413$ on the same six dimensions. This is the central
tradeoff: a shared summary is cheaper and operationally simpler, but it can
dilute dimensions whose evidence is sparse or institutionally specific.

The leaf-size sweep is also informative. In the fixed per-dimension pipeline,
macro Pearson is $0.815$ at $64\mathrm{K}$ characters, $0.824$ at
$32\mathrm{K}$, $0.824$ at $16\mathrm{K}$, $0.829$ at $8\mathrm{K}$,
$0.842$ at $4\mathrm{K}$, $0.827$ at $2\mathrm{K}$, and $0.827$ at
$1\mathrm{K}$. Economic policy stays between $0.916$ and $0.939$ across
those rows. Decentralization is lower but has a meaningful small-leaf signal:
$0.543$ at $8\mathrm{K}$, $0.580$ at $4\mathrm{K}$, and $0.584$ at
$2\mathrm{K}$. This is the local-law intuition in empirical form. Leaf
states can become small enough to inspect while the root measurement remains
visible.

\begin{figure}[!htbp]
  \centering
  \includegraphics[width=\linewidth]{manifesto_fg_combined_ladder_f1g0_f1g1.pdf}
  \caption{\textbf{First f/g step in the joint all-six ladder, by dimension.}
  External Pearson is measured on the $48$-document held-out split; each
  panel uses the held-out manifestos with expert labels for that dimension.
  The joint summary is shared across all six dimensions. The green curve is
  the best available $f^xg^y$ cell at each leaf for that dimension.
  Horizontal lines mark Benoit's proprietary LLM ensemble, the best
  open-weight Benoit reference row, and the split-expert reference for the
  same dimension.}
  \label{fig:v6-manifesto-combined-ladder}
\end{figure}

Figure \ref{fig:v6-manifesto-combined-ladder} gives the corresponding
joint all-six f/g ladder. The $f^1g^0$ macro row is $0.776$, $0.789$,
$0.773$, $0.788$, $0.783$, and $0.789$ at $256$, $512$, $1{,}024$,
$2{,}048$, $4{,}096$, and $8{,}096$ tokens. The first summarizer step,
$f^1g^1$, gives $0.764$, $0.774$, $0.803$, $0.777$, $0.785$, and
$0.789$. The best common first-step row is therefore $r=0.803$ at
$1{,}024$ tokens: above the best open-weight macro reference $0.793$ and
below Benoit's proprietary ensemble macro $0.817$ and the split-expert macro
$0.873$. If each dimension is allowed to take its best available $f^xg^y$
cell at the same leaf, the macro row is $0.782$, $0.791$, $0.806$, $0.793$,
$0.801$, and $0.809$. The dimension panels show why the macro is not the full
story. Economic, immigration, and European integration remain close to the
Benoit reference rows across the leaf sweep; environment is usable but below
the open-weight reference; decentralization is the main point of stress even
when later cells help at $8{,}096$ tokens.

For Environment, a root-only cold-start control clarifies the source of the
gap. We score the stored full-document Environment summaries with Gemma-4 and
the raw Benoit prompt; this is a single root-summary control, not a leaf-size
sweep. It reaches $r=0.784$ on the $48$-document held-out split, essentially
tied with the first ladder row at $256$ tokens. The shortfall is therefore
already present before learned tree updates; the ladder mostly preserves that
signal rather than degrading it.

\begin{table}[t]
  \centering
  \begin{tabular}{lrrrr}
    \toprule
    Dimension & Best $r$ & Leaf tokens & Stage & $n$ \\
    \midrule
    Economic & $0.886$ & $4{,}096$ & $f^2g^1$ & $47$ \\
    Social & $0.886$ & $1{,}024$ & $f^2g^2$ & $43$ \\
    Immigration & $0.938$ & $4{,}096$ & $f^3g^3$ & $29$ \\
    European integration & $0.965$ & $2{,}048$ & $f^1g^1$ & $35$ \\
    Environment & $0.795$ & $1{,}024$ & $f^1g^1$ & $42$ \\
    Decentralization & $0.461$ & $8{,}096$ & $f^1g^1$ & $42$ \\
    \bottomrule
  \end{tabular}
  \caption{\textbf{Best dimension-specific cells in the current joint
  all-six f/g optimization.} One summary is trained to preserve evidence for
  all six policy dimensions. These are best cells by dimension; a common
  deployed row would choose one leaf size and one stage.}
  \label{tab:v6-joint-dim-best}
\end{table}

The joint all-six f/g optimization makes the tension more explicit. A single
summary is asked to preserve evidence for all six scales at once. The best
dimension-specific cells in Table \ref{tab:v6-joint-dim-best} are higher
than any single deployed row for economic, social, immigration, and European
integration. Environment remains usable at $r=0.795$. Decentralization is
the outlier at $0.461$. Comparing this with the single-target
decentralization ladder and the fixed per-dimension sweep points to the
failure mode: decentralization evidence is easy for a multi-purpose summary
to lose, and optimizing the common summary can improve teacher agreement
while leaving the expert target mostly unmoved.

\subsection{Single-Dimension Ladders}\label{sec:min-manifesto-single-dim}

We first separate the single-target evidence. Economic policy is the complete
replication lane: the teacher, expert target, and tree representation agree
strongly, and the Benoit-initialized ladder is populated at every plotted leaf
size. Decentralization is the current stress lane. It is still a partial
single-target ladder, but it is the useful counterexample because Benoit's own
ensemble reaches only $r=0.490$ on this dimension, the best open-weight
reference reaches $0.450$, and the split-expert reference is
$0.780$.

\begin{figure}[!htbp]
  \centering
  \includegraphics[width=0.90\linewidth]{manifesto_fg_econ_decent_f1g0_f1g1.pdf}
  \caption{\textbf{Single-dimension f/g ladders for economic policy and
  decentralization.}
  External Pearson is measured against Benoit expert means on the $48$-document
  held-out split. The economic panel uses the complete Benoit-initialized
  economic lane; because $g^0$ is Benoit's summary in that lane, the plotted
  initial row is $f^1g^0$. The decentralization panel shows the current
  single-target lane: $f^1g^0$ is available through $4{,}096$ tokens and
  $f^1g^1$ through $2{,}048$ tokens. In both panels, the green curve is the
  best available $f^xg^y$ cell at each leaf. Horizontal lines mark Benoit's
  proprietary LLM ensemble, the best open-weight reference value, and the
  split-expert reference for the same dimension.}
  \label{fig:v6-manifesto-econ-decent-ladder}
\end{figure}

Figure \ref{fig:v6-manifesto-econ-decent-ladder} compares the tuned query on
the initial summaries, $f^1g^0$, with the first optimized summarizer step
where that step is available, $f^1g^1$. On economic policy, the complete
Benoit-initialized $f^1g^0$ row gives $r=0.824$, $0.861$, $0.885$, $0.874$,
$0.865$, and $0.866$ at $256$, $512$, $1{,}024$, $2{,}048$, $4{,}096$, and
$8{,}192$ token leaves. The best $f^xg^y$ cell at each economic leaf is
$0.830$, $0.861$, $0.885$, $0.886$, $0.886$, and $0.879$. Thus the complete
economic lane stays in the same range as Benoit's proprietary ensemble
($0.870$), the best open-weight reference ($0.860$), and the split-expert
reference ($0.880$), even as the leaf size changes by a factor of $32$. The
separate raw-init economic branch is not the main plotted lane because it is
only complete through $1{,}024$ tokens; it does contain the current best
single economic cell, $r=0.909$ at $256$ tokens.

Decentralization behaves differently. At $256$ tokens, the first summarizer
step raises external Pearson from $0.534$ to $0.557$. At $512$ tokens it moves
from $0.476$ to $0.449$, and at $1{,}024$ tokens it moves from $0.510$ to
$0.435$. At $2{,}048$ tokens it moves from $0.479$ to $0.452$. Later ladder
cells recover part of the loss at $1{,}024$ tokens: $f^2g^2$ reaches
$r=0.536$. The best available decentralization cells by leaf are $0.557$,
$0.476$, $0.536$, $0.479$, and $0.485$ at $256$, $512$, $1{,}024$,
$2{,}048$, and $4{,}096$ tokens. These single-target cells improve on the
current joint all-six optimized run on decentralization, whose best
dimension-specific cell is $0.461$, but they remain well below the
split-expert reference of $0.780$. This is the contrast we want:
decentralization separates compression quality from teacher/expert alignment.
The internal teacher correlation stays high while external expert correlation
moves much less, so the ladder shows where the model is learning the teacher
distribution while only partly recovering the expert target.

\paragraph{Scope of the single-target ladders.}
The same-split individual f/g ladder is complete for economic policy and
partially populated for decentralization. The six-dimension headline in
Table~\ref{tab:min-benoit-headline} is the complete root-observed benchmark;
the single-target ladders are diagnostic slices showing how alternating
query and summarizer optimization behaves on one strong dimension and one
stress dimension.

\subsection{Audit-Gap Diagnostics}\label{sec:min-manifesto-audit-gap}

The available local diagnostics come from teacher traces and local residual
checks. They are useful for seeing where the learned state drifts, but they
should not be read as the final design-based human audit from
Section~\ref{sec:v7-neural-operator-route}. In the notation of
Section~\ref{sec:v7-tree-comparisons}, they are empirical local-law
measurements for the fitted LLM pipeline, while the root-observed benchmark
above remains the external validation target.

\begin{figure}[!htbp]
  \centering
  \includegraphics[width=\linewidth]{manifesto_fg_combined_audit_gap.pdf}
  \caption{\textbf{Joint all-six local diagnostic gaps.}
  The figure reports available teacher-trace/local residual diagnostics for
  the joint all-six ladder. These gaps diagnose where the learned state and
  query drift across leaf and merge calls; they are not a substitute for a
  future human node-level audit with logged propensities.}
  \label{fig:v6-manifesto-combined-audit-gap}
\end{figure}

\begin{figure}[!htbp]
  \centering
  \includegraphics[width=0.90\linewidth]{manifesto_singledim_per_dim_live_audit_gap.pdf}
  \caption{\textbf{Single-dimension local diagnostic gaps.}
  Economic policy and decentralization illustrate the same distinction as the
  root ladders: one dimension largely preserves the external target, while the
  other exposes teacher/expert and summary-capacity stress.}
  \label{fig:v6-manifesto-single-audit-gap}
\end{figure}

\clearpage

\subsection{Takeaways}\label{sec:min-manifesto-takeaways}

The Manifesto result should be read as an auditable-compression result. The
per-dimension tree reaches the published frontier range at the root, and the
combined tree remains close on five of six dimensions, while every run
produces a finite population of leaves, summaries, and merge states. That is
the structure flat scoring and concat-no-merge controls lack. They are useful
controls because they show that long-context scoring can be strong; the tree
adds local units that can be trained, read, sampled, and audited.

The next measurement layer is human node supervision. The current results use
root expert means for evaluation and teacher traces for node training. The
Manifesto Project quasi-sentence codings can be aggregated inside sampled
leaf or merge spans, turning the same tree population into a source of local
C1 and C3 labels. Root labels continue to calibrate the document-level expert
target; quasi-sentence aggregates supply local evidence labels; sampled nodes
then feed the audit certificate. The practical question has shifted from
whether the root score can survive tree compression to how much human node
labeling is needed to certify that the surviving score is carried by locally
faithful states.

The limitations are the same distinctions used throughout
Section~\ref{sec:v7-tree-comparisons}. Exact classical parity is a
state-level claim; learned LLM compression is an empirical approximation;
dimensions such as decentralization expose structural or measurement gaps;
and audit uncertainty belongs to the sampling design. The paper therefore
claims a complete root-observed LLM result here, plus a concrete route to the
stronger node-level certificate.
